\section{\texttt{gdt\_init()}：填充三个描述符}

\codefile{src/kernel/arch/gdt.c}
\begin{lstlisting}[style=cstyle]
void gdt_init(void) {
    // 描述符 0：Null 描述符（x86 硬性要求）
    //   所有字段都是 0。
    //   如果任何段寄存器的 selector = 0x00，CPU 知道它是"空的"。
    //   在某些场景（如加载 LDT 寄存器）CPU 会检查是否为 null selector。
    gdt_set_entry(0, 0, 0, 0, 0);

    // 描述符 1：内核代码段（selector = 0x08）
    gdt_set_entry(1,
        0x00000000,   // base  = 0（平坦模型）
        0x000FFFFF,   // limit = 0xFFFFF（配合 G=1 → 4GiB）
        0x9A,         // access: P=1, DPL=0, S=1, Ex=1, DC=0, RW=1, Ac=0
        0xC0          // flags:  G=1, D/B=1, L=0, AVL=0
    );

    // 描述符 2：内核数据段（selector = 0x10）
    gdt_set_entry(2,
        0x00000000,   // base  = 0（平坦模型）
        0x000FFFFF,   // limit = 0xFFFFF
        0x92,         // access: P=1, DPL=0, S=1, Ex=0, DC=0, RW=1, Ac=0
        0xC0          // flags:  G=1, D/B=1, L=0, AVL=0
    );

    gdtr.limit = (uint16_t)(sizeof(gdt) - 1);  // 23
    gdtr.base  = (uint32_t)(uintptr_t)&gdt[0];

    gdt_flush(&gdtr);
}
\end{lstlisting}

\begin{figure}[H]
  \centering
  % figures/ch03/fig08-gdt-3-24.tex — auto-extracted from chapter source
\begin{tikzpicture}[node distance=0cm]
  % GDT array in memory
  \node[memcell, minimum width=9cm, minimum height=0.8cm, fill=gray!10] (e0) {
    00 00 00 00 00 00 00 00
  };
  \node[left=0.5cm of e0, font=\small\bfseries] {gdt[0] (Null)};

  \node[memcell, minimum width=9cm, minimum height=0.8cm, fill=green!10, above=0cm of e0] (e1) {
    FF FF 00 00 00 9A CF 00
  };
  \node[left=0.5cm of e1, font=\small\bfseries] {gdt[1] (Code)};

  \node[memcell, minimum width=9cm, minimum height=0.8cm, fill=blue!10, above=0cm of e1] (e2) {
    FF FF 00 00 00 92 CF 00
  };
  \node[left=0.5cm of e2, font=\small\bfseries] {gdt[2] (Data)};

  % GDTR arrow
  \node[left=3.5cm of e0, font=\small] (gdtr) {GDTR};
  \draw[pointer] (gdtr.east) -- (e0.west);
  \node[below=0.1cm of gdtr, font=\scriptsize\ttfamily] {limit=23, base=\&gdt};
\end{tikzpicture}

  \caption{GDT 初始化后的内存布局（3 个描述符 = 24 字节）}
\end{figure}

\begin{designbox}
为什么 GDT 第 0 项必须是空描述符？

Intel 规范规定 selector=0 是 null selector。作用有二：
\begin{enumerate}
  \item \textbf{防御性检查}：如果代码不小心用了未初始化的段寄存器（值为 0），
    CPU 在访问内存时会检测到 null selector 并触发 \#GP，而非悄悄使用错误的段。
  \item \textbf{LDT 相关}：加载 LDT 寄存器时，null selector 表示"不使用 LDT"。
\end{enumerate}
\end{designbox}

\begin{thinkbox}
将来做用户态（Ring 3）时，我们需要添加 2 个新描述符：
\begin{itemize}
  \item 用户代码段（DPL=3，access=\hex{FA}）
  \item 用户数据段（DPL=3，access=\hex{F2}）
\end{itemize}
它们的 selector 分别是多少？（提示：index 3 和 4，各乘以 8。）
\end{thinkbox}

% ============================================================================

